\section{The complete coefficient span of a one-particle digit family}
\label{sec:one-particle-digits}

The sparse-digit construction has a particularly explicit family
when both symbol sizes are one. We determine its entire space of
coefficient tensors and prove a rank obstruction for every scalar
linear reconstruction in that space. This restriction concerns the
specified family; it does not apply to larger symbol sizes or to
arbitrary crossing decompositions.

Work in characteristic two. Fix $q=2^p\ge4$ and $h\ge1$, and put
\[
 m=2^h,\quad H=m-1,\quad
 S=\sum_{j=0}^{h-1}q^j,\quad e=2(q-1)S+2=2q^h.
\]
Take $n_1=n_2=1$ in Section~\ref{sec:deformed-crossing}, with the
full symbol class of energy $e$. It is precisely
\[
 \Sigma=\{(\{i\},\{H\mathbin{\mathrm{xor}}i\}):0\le i<m\}.
\]
Indeed, $E_q(i)+E_q(j)=S$ has no carries in radix $q\ge4$:
each digit of either summand is zero or one, and the target digits
are all one in the first $h$ positions. Thus the two bit strings
are complementary. Label the parties by $a,b,c,d\in\{0,\ldots,m-1\}$
before exchanging the two sets of the symbols at $c,d$.
Write
\[
 C_h(a,b,c,d)=[a=H\mathbin{\mathrm{xor}}c]
                         [b=H\mathbin{\mathrm{xor}}d].
\]
This is the crossing tensor up to those local permutations.

\begin{proposition}[Three surviving determinant terms]
\label{prop:one-particle-formula}
The determinant in Theorem~\ref{thm:curve-crossing} satisfies
\begin{align}
 \Psi={}&X^{2e}C_h
   +f^2y^{E_q(a\mathbin{\mathrm{xor}}d)+
                    E_q(b\mathbin{\mathrm{xor}}c)}\notag\\
 &\quad+f^2X^4y^{4S-E_q(a\mathbin{\mathrm{xor}}b)-
                           E_q(c\mathbin{\mathrm{xor}}d)}.
 \label{eq:one-particle-formula}
\end{align}
This identity holds with $f$ an independent indeterminate,
before imposing its curve equation.
\end{proposition}

\begin{proof}
For labels $i,j$ set
\[
 H(i,j)=y^{E_q(i\mathbin{\&}(H\mathbin{\mathrm{xor}}j))},
 \qquad
 L(i,j)=Xy^{E_q(i)+E_q(H\mathbin{\mathrm{xor}}j)-
                   E_q(i\mathbin{\&}(H\mathbin{\mathrm{xor}}j))}.
\]
These are respectively the $C$ entry and the $D_iB_{ij}$ entry
of the determinant, with the column index complemented.
For any two labels the ratio of a nonzero matrix entry to its
transposed entry is $y^{E_q(i)-E_q(j)}$; the assertion also holds
for the Kronecker-delta entries on their support.
Multiplying the row with label $i$ by $y^{-E_q(i)}$ therefore
makes the matrix symmetric with zero diagonal.
In a four-by-four symmetric zero-diagonal matrix in characteristic
two, the six four-cycle terms of the determinant cancel in reverse
pairs. Only its three products of two transpositions remain.
Undoing the row scalings gives the same conclusion for $\Psi$.

The transpositions $(a,c)(b,d)$ give $X^{2e}C_h$.
For the other terms, digitwise calculation gives
\[
 H(i,j)H(j,i)=y^{E_q(i\mathbin{\mathrm{xor}}j)},\qquad
 L(i,j)L(j,i)=X^2y^{2S-E_q(i\mathbin{\mathrm{xor}}j)}.
\]
The transpositions $(a,d)(b,c)$ and $(a,b)(c,d)$ thus give the
second and third terms in \eqref{eq:one-particle-formula}.
The calculation in the function field is an identity of polynomials
in $X,f$, so no nonzero-coordinate restriction remains.
\end{proof}

For binary labels define $\delta(i,j)=[i=j]$ and
$\eta(i,j)=[i\ne j]$. Introduce the four-party tensors
\[
\begin{array}{lll}
 A_0=\delta(a,d)\delta(b,c),&
 A_1=\eta(a,d)\delta(b,c)+\delta(a,d)\eta(b,c),&
 A_2=\eta(a,d)\eta(b,c),\\
 B_0=\delta(a,b)\delta(c,d),&
 B_1=\eta(a,b)\delta(c,d)+\delta(a,b)\eta(c,d),&
 B_2=\eta(a,b)\eta(c,d).
\end{array}
\]
Let $\mathcal A=\operatorname{span}\{A_0,A_1,A_2\}$ and
$\mathcal B=\operatorname{span}\{B_0,B_1,B_2\}$, and put
$C=\eta(a,c)\eta(b,d)$. Products group the $h$ bits at each party,
so $C_h=C^{\otimes h}$.

\begin{proposition}[The exact coefficient space]
\label{prop:one-particle-span}
The space spanned by all coefficient tensors of the function $\Psi$
on the Fermat curve is
\[
 \mathcal S_h=\mathcal A^{\otimes h}+
                  \mathcal B^{\otimes h}+\operatorname{span}\{C_h\}.
\]
Its dimension is $2\cdot3^h-2^h+1$ when $h\ge2$, and four when
$h=1$. Every tensor obtained by a scalar linear functional on the
coordinate-function space of $\Psi$ belongs to this space, including
all jets at either distinguished point after the stated scalar
normalizations and reflection.
\end{proposition}

\begin{proof}
The exponents $E_q(a\mathbin{\mathrm{xor}}d)+
E_q(b\mathbin{\mathrm{xor}}c)$ have independent radix-$q$ digits
in $\{0,1,2\}$. Their coefficient tensors in the second term of
\eqref{eq:one-particle-formula} are exactly the products of the
$A_0,A_1,A_2$. The third term gives exactly the products of the
$B_0,B_1,B_2$. The scalar monomials within each family are distinct.
A collision between the two families would require
$2(q-1)$ to divide four, which is impossible for odd $q-1\ge3$.
Finally $f^2=Y^4$ is not rational in $X$: the extension
$\kk(X,Y)/\kk(X)$ has degree $q-1$, by Eisenstein at $X-1$,
and $\gcd(4,q-1)=1$ means that rationality of $Y^4$ would imply
rationality of $Y$ itself. Thus the rational function $X^{2e}$
is independent of all the scalar functions $f^2X^j$ above.
These facts prove the asserted coefficient space.

Let $O$ be the indicator of odd total bit parity and $E$ the
indicator of even total bit parity. Direct evaluation gives
\[
 A_1=B_1=O,\qquad A_0+A_2=B_0+B_2=E,\qquad C=A_0+B_0.
\]
The tensors $E,O,A_0,B_0$ are linearly independent: the coordinates
$0101,0001,0110,0011$, in that order, isolate the necessary
coefficients. Hence
$\mathcal A\cap\mathcal B=\operatorname{span}\{E,O\}$.
Choose this common basis and extend it by $A_0$ on one side and
$B_0$ on the other. The resulting tensor-product basis shows
\[
 \mathcal A^{\otimes h}\cap\mathcal B^{\otimes h}
       =\operatorname{span}\{E,O\}^{\otimes h}.
\]
For $h\ge2$, the expansion of $(A_0+B_0)^{\otimes h}$ contains
mixed $A_0,B_0$ basis words, so $C_h$ is outside the sum of the
two spaces. The dimensions follow. For $h=1$, $C$ already belongs
to their four-dimensional sum.

Multiplying all scalar functions by a fixed nonzero function and
applying a scalar functional cannot enlarge the tensor span.
Reflection exchanges $b,d$, which exchanges $\mathcal A$ with
$\mathcal B$ and fixes $C$. It consequently preserves $\mathcal S_h$.
This covers the normalized jets at infinity as well as those at zero.
\end{proof}

\begin{theorem}[A rank obstruction for every linear reconstruction]
\label{thm:one-particle-rank-barrier}
Let $h\ge2$ and let $\alpha\ne0$ belong to any field of
characteristic two. Every tensor
\[
 T\in\alpha C_h+\mathcal A^{\otimes h}+\mathcal B^{\otimes h}
\]
has rank at least $(4^h-2^h)/2$ across each of the two adjacent cuts.
Consequently, if $C_h=T_1+T_2$ with $T_1,T_2\in\mathcal S_h$, then
\[
 \operatorname{rk}_{ad\mid bc}T_1+
 \operatorname{rk}_{ab\mid cd}T_2\ge\frac{4^h-2^h}{2}.
\]
The minimum such rank sum has exponent exactly two in the wire
dimension $2^h$ as $h\to\infty$.
\end{theorem}

\begin{proof}
First flatten across $ad\mid bc$. In each four-dimensional pair
space use the basis
\[
 x=e_{00}+e_{11},\quad y=e_{01}+e_{10},\quad
 z=e_{00},\quad w=e_{01}.
\]
For a two-by-two matrix $R$ write
\[
 D(R)=\begin{pmatrix}0&R\\0&0\end{pmatrix}
\]
relative to the decomposition
$\operatorname{span}(x,y)\oplus\operatorname{span}(z,w)$.
Let $P=\operatorname{diag}(1,0)$ and
$J=\left(\begin{smallmatrix}0&1\\1&0\end{smallmatrix}\right)$.
The displayed tensors have the following exact matrix forms:
\[
 \begin{array}{lll}
 A_0=D(P),&A_1=D(J),&A_2=D(I-P),\\
 B_0=I,&B_1=D(J),&B_2=I+D(I),\qquad C=I+D(P).
 \end{array}
\]
These identities follow by applying each operator to $x,y,z,w$.

Give a product of $D$ factors degree equal to the number of such
factors; identity factors have degree zero. Since $D(R)D(R')=0$,
all positive-degree operators strictly raise the number of factors
in $\operatorname{span}(x,y)$. If the degree-zero coefficient of
$T$ is nonzero, $T$ is a nonzero scalar identity plus a nilpotent
operator, hence has full rank $4^h$.

Otherwise, since $h\ge2$, the summand from $\mathcal A^{\otimes h}$
does not affect degree one. The degree-one part of $T$ is
\[
 L=\sum_{i=1}^h I^{\otimes(i-1)}\otimes D(R_i)
                       \otimes I^{\otimes(h-i)},\qquad
 R_i=\alpha P+a_iI+b_iJ.
\]
In characteristic two, $\operatorname{tr}R_i=\alpha\ne0$.
Thus $\rho_i=\operatorname{rk}D(R_i)=\operatorname{rk}R_i\ge1$.
Conjugate each pair factor by
$\operatorname{diag}(tI_2,I_2)$, and divide the full operator by $t$.
Over the rational function field this preserves rank, and its limit
at $t=0$ is $L$. Vanishing of matrix minors gives
$\operatorname{rk}T\ge\operatorname{rk}L$.

Each square-zero operator $D(R_i)$ is a direct sum of $\rho_i$
Jordan blocks of size two and $4-2\rho_i$ zero blocks of size one.
On a tensor product of these blocks, $L$ is the sum of their
individual differentials. If at least one factor is a size-two
block, choose its inverse-direction map $H$. Then
$LH+HL=I$: on that factor this is the elementary Jordan-block
identity, and the cross terms cancel because the characteristic
is two. Hence the kernel and image of $L$ coincide on that summand,
and $L$ has half its dimension as rank. The only remaining summands
are products of zero blocks, with total dimension
$\prod_i(4-2\rho_i)\le2^h$. Therefore
\[
 \operatorname{rk}L=
 \frac{4^h-\prod_i(4-2\rho_i)}2
 \ge\frac{4^h-2^h}2.
\]
This proves the bound for the first cut. Exchanging $b,d$ interchanges
$\mathcal A,\mathcal B$, preserves $C_h$, and interchanges the cuts,
so it proves the other bound as well.

By Proposition~\ref{prop:one-particle-span}, $C_h$ is independent
of $\mathcal A^{\otimes h}+\mathcal B^{\otimes h}$.
In a decomposition $C_h=T_1+T_2$, the two coefficients of $C_h$
therefore sum to one. At least one is nonzero, and its appropriate
cut rank has the preceding lower bound. The trivial decomposition
$C_h+0$ has rank sum $4^h$. These two bounds have logarithmic
exponent two in $2^h$, proving the final assertion.
\end{proof}

This theorem covers arbitrary scalar combinations of the entire
coefficient space, including combinations chosen separately for each
$h$. It therefore includes a reconstruction based on fewer scalar
functions than the ambient Riemann--Roch space. It does not cover
larger symbol sizes, new kernels, or local tensor transformations
that leave this coefficient space.

The One-particle digit verifier in the workbench checks
\eqref{eq:one-particle-formula} on 4640 entries for the stated small
radices and dimensions. It also enumerates all $16384$ tensors in
$C_2+\mathcal A^{\otimes2}+\mathcal B^{\otimes2}$ over $\F_2$:
the minimum first-cut rank is six, attaining the theorem's bound.
These finite checks support the proof and do not replace its
all-$h$ argument.
